% ═══════════════════════════════════════════════════════════════════════════
%  EXPERIMENT I — Baseline single-run comparison  (UPDATED with real data)
%  Requires: \usepackage{booktabs, siunitx, amsmath, amssymb}
% ═══════════════════════════════════════════════════════════════════════════

%  NOTE: The unified file "implementation_and_pareto.tex" contains
%  the canonical version of this section.  This file is the standalone
%  backup with the same corrected numbers.

Table~\ref{tab:exp1} summarises the steady-state tracking performance
of DQ-MPCC and the quaternion-based MPCC baseline over a single
\SI{100}{m} lap of the Lissajous reference path.
Both controllers use their respective Optuna-tuned weights,
identical constraints ($T_{\max} = 10g$, $\tau_{\max} = \SI{0.5}{Nm}$,
$v_{\theta,\max} = \SI{15}{m/s}$), prediction horizon
($T_p = \SI{0.3}{s}$, $N = 30$ nodes), and control frequency
(\SI{100}{Hz}); the only difference is the state-space
representation.

\begin{table}[t]
  \centering
  \caption{Experiment~I — single-run comparison
           ($S_{\max}=\SI{100}{m}$, $v_{\theta,\max}=\SI{15}{m/s}$).}
  \label{tab:exp1}
  \setlength{\tabcolsep}{4pt}
  \begin{tabular}{lSS}
    \toprule
    {Metric} & {DQ-MPCC} & {MPCC Baseline} \\
    \midrule
    {RMSE$_{\text{pos}}$ [m]}       & 0.972 & 0.688 \\
    {RMSE$_{\text{cont}}$ [m]}      & 0.901 & 0.612 \\
    {RMSE$_{\text{lag}}$ [m]}       & 0.365 & 0.314 \\
    {Mean $|\mathbf{e}_p|$ [m]}     & 0.790 & 0.430 \\
    {Max $|\mathbf{e}_p|$ [m]}      & 3.019 & 3.000 \\
    {$\bar v_\theta$ [m/s]}         & 11.22 & 12.36 \\
    {$t_{\text{lap}}$ [s]}          & 8.91  & 8.10  \\
    {Effort $\sum\|\mathbf{u}\|^2$} & 1258.6 & 836.7 \\
    {$\bar t_{\text{solver}}$ [ms]} & 5.40  & 3.23  \\
    {$t_{\text{solver}}^{\max}$ [ms]} & 9.57  & 5.42  \\
    {$N_{\text{steps}}$}            & 891   & 810   \\
    \bottomrule
  \end{tabular}
\end{table}

In this nominal scenario, the baseline achieves lower position
RMSE (\SI{0.688}{m} vs.\ \SI{0.972}{m}) and completes the lap
$10\%$ faster ($t_{\text{lap}} = \SI{8.10}{s}$ vs.\ \SI{8.91}{s}).
The baseline's mean progress speed is $\bar v_\theta = \SI{12.36}{m/s}$
versus DQ-MPCC's \SI{11.22}{m/s}, meaning it traverses the path
more aggressively.
The DQ-MPCC solver requires \SI{5.40}{ms} per step versus
\SI{3.23}{ms} for the baseline, attributable to the larger
state dimension ($n_x = 15$ vs.~$14$) and the $\mathrm{SE}(3)$
logarithm in the cost; both remain well within the \SI{10}{ms}
real-time budget.

At first glance, these results appear to favour the baseline.
However, comparing RMSE at a \emph{single} speed and initial
condition conflates tracking accuracy with the controller's
aggressiveness profile.
To disentangle speed and accuracy and to evaluate statistical
robustness, we turn to Experiment~II.
